\section{Related work}
\label{sec:related}

\paragraph{Honesty versus accuracy.}
MASK~\citep{ren2025mask} separates \emph{honesty} (does the statement match the model's elicited belief?) from \emph{accuracy} (does the belief match the truth?) by eliciting beliefs neutrally and then applying pressure prompts. It is a behavioural benchmark and does not analyse detectors. \citet{pacchiardi2023catch} and \citet{evans2021truthful} use the same definition of a lie as a false statement made despite knowing the truth, and \citet{campbell2023localizing} open by noting that it is often unclear whether a false output reflects missing knowledge or dishonesty, but then study instructed lying only. Our design takes the MASK decomposition, applies it to matched questions for a single model, and asks what white-box detectors do in each cell.

\paragraph{Deception probes and lie-detection benchmarks.}
\citet{goldowskydill2025detecting} train linear probes on Llama-3.3-70B from contrastive honest/deceptive instructions~\citep{zou2023representation} and from role-play scenarios, report AUROCs of 0.96--0.999 on insider-trading and sandbagging settings, and note that the probes also fire on honest text that is merely about deception. Liars' Bench~\citep{kretschmar2025liarsbench} collects 72,863 lies and honest responses from four open-weight models and finds that detectors fail systematically on some lie types; its negatives are honest responses, not honest errors. \citet{natarajan2026oneprobe} show that deception types have partly distinct representations and that prompt choice explains most of the variance in probe performance; \citet{moustafa2026beyond} study how depth, probe expressivity and lie typology (fabrication, omission, exaggeration) affect probes; \citet{luikham2026asymmetries} find asymmetric transfer between probes for spontaneous and instructed deception; \citet{thormann2026probing} show that truth probes miss deception that uses no false statement; \citet{parrack2025benchmarking} and \citet{boxo2025caught} benchmark deception probes on further models. \citet{cooney2026didyoulie} argue that lie-detector evaluations need settings where the model verifiably believes the opposite of what it says, and \citet{hopkins2026finetuned} elicit on-policy lies with neutral belief elicitation followed by pressure and find that fine-tuned ``did you lie?'' detectors do not generalise across lie categories; they also report that many labelled lies may reflect confusion rather than knowing deception. All of these compare lies with honest (mostly correct) responses or one lie type with another. None of them, as far as we can tell from the papers' abstracts and stated designs, crosses a knowledge label with an incentive manipulation on matched questions and asks whether a detector separates a lie from an honest error.

\paragraph{Truth probes.}
\citet{azaria2023internal}, \citet{burns2022discovering}, \citet{marks2023geometry} and \citet{burger2024truth} find linear directions separating true from false statements. \citet{levinstein2023still} argue that such probes do not generalise and are not lie detectors. A truth probe is, by construction, a falsehood detector; whether it can separate a lie from a hallucination, both false, is the question we test.

\paragraph{Hallucination detection and knowledge.}
Uncertainty-based detectors (semantic entropy~\citep{kuhn2023semantic,farquhar2024detecting}; its probe variant~\citep{kossen2024semantic}) and error-detection probes~\citep{orgad2024llms} predict whether an answer is wrong. \citet{cheang2025really} show that hidden states mainly reflect whether the model is recalling parametric knowledge rather than whether the output is true, so that some hallucinations look like correct answers internally. \citet{simhi2024distinguishing,simhi2025trust} split hallucinations by whether the model has the knowledge, \citet{slobodkin2023curious} find that unanswerability is encoded even when the model hallucinates, and \citet{gekhman2025insideout} operationalise ``knowing'' through internal ranking of answers. These works vary knowledge but not incentive. We use a hallucination probe trained in the standard way and ask how it treats lies.

\paragraph{Incentive-driven misreporting.}
Sandbagging~\citep{vanderweij2024sandbagging}, sycophancy~\citep{sharma2023sycophancy} and strategic deception under pressure~\citep{scheurer2023strategically,park2024deception} are the behaviours that lie detectors are ultimately meant to catch. Our pressure prompts are simple single-turn versions of these (a game incentive, a goal conflict, and a deployment threat) plus an instructed and a role-play condition for comparison.
